\chapter{VLM 怎样调用技能并形成闭环}
\label{chap:runtime}

这套系统把“理解画面与指令，选择下一步”和“让关节持续运动并保持平衡”放在不同层。上层输出一个技能名称及受限时长；下层用已有 ONNX 策略产生连续关节目标。训练改变的是上层选择方式，现有走路、平衡和左右踢球模型在这次微调中不更新。

\section{还没有机器人微调时，模型为什么知道技能库}

模型并不会自动扫描本机目录，也不会仅凭“你是一只小鸭子”知道有哪些动作。应用程序每次发送完整的 \texttt{SYSTEM\_PROMPT}，其中明确列出九个技能、机身坐标、方向含义、持续时间范围、两台相机的用途、脚边目标点以及回复格式。它还包含“走一步后结束”“看不见目标先转身寻找”等接口示例。这是人写的使用说明，属于推理时上下文；没有因发送提示词而更新模型权重。

预训练的指令模型已有读取图像和语言、遵循指令、生成文本的基础，因此可以尝试在这个说明下选择技能。它是否选得正确、能否持续执行完整任务，仍需测试。微调则把大量“本机画面＋原指令＋必要状态和历史 $\rightarrow$ 正确技能及参数”变成训练监督，让权重更适应这套机器人和任务。

因此至少要区分三层“已经训练”：Qwen 的通用预训练／指令训练、已有 ONNX 动作策略的训练、本次机器人技能调用的微调。代码里的 \texttt{is\_trained\_vlm=True} 只说明正在调用一个有预训练权重的模型；同处的 \texttt{robot\_finetuned=False} 明确没有加载本项目微调结果。不能把前者读成“本机微调已经完成”。

\subsection{目标颜色究竟从哪里来}

以“别踢蓝球，把红球踢出去”为例，原话直接进入 user 输入。模型需要理解要操作的是红球，再从双图和辅助检测寻找对应目标；当前技能层不调用规则解析器先替它决定红色。提示词只告诉它合法颜色值有 red、yellow、blue，没有替当前指令填写答案。

模型输出的 \texttt{target\_color} 是它本次声明的目标，不保证一定理解正确。guard 检查的是这个声明颜色的球是否可见、是否在脚边，不能代替语言理解。如果模型把红球任务说成蓝球，而蓝球又恰好符合脚边条件，仅通过 guard 也不等于完成用户任务；独立测试的 \texttt{expected\_color} 才会按预先指定的目标验收。该测试参数不送给模型。

历史 \texttt{selected\_color} 可以来自以前实际选择过的任务；用户当前明确指定另一颜色时应以新指令为准。普通前进、后退和问答的 \texttt{target\_color} 应为 null，不能因为场景里有红球就擅自生成红球任务。

\section{九个技能和五字段 JSON}

输出格式由项目的 \texttt{PROPOSAL\_SCHEMA} 与 \texttt{parse\_proposal} 定义，不是 Qwen 的天然固定输出。通用模型本来可以生成普通对话；本项目通过提示词、服务端结构约束和客户端严格解析把它限定为下表协议。

\begin{longtable}{P{3.8cm}P{6.1cm}P{4.6cm}}
\hline
字段 & 含义 & 合法值／约束\\
\hline
\endhead
\texttt{type} & 本次要执行技能、回答、澄清还是结束。 & skill、answer、clarify、finish。\\
\texttt{skill} & 要调用的已注册动作。 & 九选一；非 skill 类型必须为 null。\\
\texttt{seconds} & 本次短动作的持续时间，单位秒。 & 0.02\,s 的整数倍；非动作类型为 0。\\
\texttt{target\_color} & 此任务要操作的球颜色。 & red、yellow、blue 或 null。\\
\texttt{reply} & 面向人的简短中文说明。 & 字符串，最多 1000 字符。\\
\hline
\end{longtable}

\begin{longtable}{P{4.3cm}P{6.3cm}P{3.9cm}}
\hline
技能名称 & 底层请求 & 上层允许时长\\
\hline
\endhead
\texttt{stand} & 零速度指令，持续平衡。 & 0.02--1.00\,s\\
\texttt{move\_forward} & 前向速度指令 $+0.2$\,m/s。 & 0.02--0.60\,s\\
\texttt{move\_backward} & 前向速度指令 $-0.2$\,m/s。 & 0.02--0.60\,s\\
\texttt{move\_left} & 横向速度指令 $+0.2$\,m/s。 & 0.02--0.60\,s\\
\texttt{move\_right} & 横向速度指令 $-0.2$\,m/s。 & 0.02--0.60\,s\\
\texttt{turn\_left} & 偏航角速度指令 $+1.0$\,rad/s。 & 0.02--1.00\,s\\
\texttt{turn\_right} & 偏航角速度指令 $-1.0$\,rad/s。 & 0.02--1.00\,s\\
\texttt{kick\_left} & 调用已有左脚踢球 ONNX。 & 固定 0.50\,s\\
\texttt{kick\_right} & 调用已有右脚踢球 ONNX。 & 固定 0.50\,s\\
\hline
\end{longtable}

速度只是给行走策略的请求，短脉冲实际位移受启动、步态和动力学影响，不保证等于“速度乘时间”。技能类别是离散的，但时长还要预测，并以 20\,ms 离散化；因此比只做九分类多一个参数选择问题。本轮采用自回归 JSON 监督微调，没有另外增加九分类动作头，也没有训练 VQ-VAE 动作码本。

例如，下面是\textbf{教学假设}的格式演示，不是模型测试输出：

\begin{Verbatim}[breaklines=true,fontsize=\small]
{"type":"skill","skill":"move_left","seconds":0.20,
 "target_color":"red","reply":"向左小幅调整，再观察红球。"}
\end{Verbatim}

解析器要求恰好这五个键，不接收 Markdown 包裹、额外字段、重复键、动作列表或关节代码。它拒绝 NaN、无穷数、错误类别、非法时间和不满足 0.02\,s 步长的参数。上层 stand 最大 1\,s，而底层执行器可执行踢后 3\,s 恢复；后者是执行协议，不能由模型绕过上层约束随意指定。

本节核对：\texttt{skill\_vlm.py}:22--98、153--182；\texttt{skills.py}:14--21。

\section{一轮真实运行经过哪些部件}

\begin{enumerate}
\item \textbf{观察。} 取得当前双图和测量状态，构造 \texttt{SkillInput}，四项内容为 \texttt{instruction}、\texttt{images}、\texttt{state}、\texttt{context}。保留原话及最近六条动作／反馈；去除任务编号，使用任务内相对时间。
\item \textbf{请求模型。} \texttt{LocalVLMBackend} 把 RGB 编成 PNG，再作为图像内容发送到本机回环 HTTP 服务。消息含 system 与 user，温度为 0，最多生成 256 tokens，并请求严格 JSON schema；模型别名是 \texttt{microduck-vlm}。
\item \textbf{解析与新观测检查。} 先严格解析 JSON，再取得更新的观测做 guard。模型推理需要时间，初始图片不能直接被当作执行瞬间的状态。
\item \textbf{执行或拒绝。} guard 通过才调用 \texttt{Runtime.execute}。拒绝时不偷偷替换成规则动作，而把拒绝原因留给下一次模型决策。连续三次拒绝则结束任务。
\item \textbf{底层控制。} 对应 ONNX 读取本体状态并连续输出关节目标，MuJoCo 推进物理。记录实际开始、结束以及 completed 状态。
\item \textbf{恢复、反馈、再看。} 非站立动作自动恢复站立，普通动作 1\,s、踢球 3\,s；然后重新观察，把本次提议及实际反馈放进 history，继续同一语言任务。
\end{enumerate}

guard 对动作要求机身倾角不超过 $25^\circ$；踢球还要求身体静止、指定颜色在当前 RGB-D 测量中存在，以及球心落在对应脚目标点的双轴 $\pm0.012$\,m 区域内。它只作接受或拒绝，不替模型决定该左移还是右移。其目标也不是覆盖全部导航危险：它没有把普通移动自动变成避障规划；相关选择仍由模型负责。

上一条记录即使叫 \texttt{previous\_action}，也要结合 \texttt{previous\_feedback} 理解。部署时其中可能保存被拒绝的提议，不能因为历史里出现过 kick 就认为真的踢了；只有成功执行反馈才能证明执行。当前成功示范主要提供 \texttt{executed\_and\_recovered} 反馈，不能宣称已经充分训练了所有拒绝后的恢复行为。

推理超时、格式错误或服务异常会留下日志并结束当前任务；程序没有静默回退到规则老师。停止请求有独立通道，推理在后台线程等待时仍可取消；被取消或超时请求的迟到回复不会再进入动作执行器。

\section{50 Hz、200 Hz 和 VLM 推理不是一个频率}

MuJoCo 每 0.005\,s 走一步，执行器 controller 也在每步更新，因此是 200\,Hz。技能策略每 0.02\,s 计算一次 ONNX 输出，即 50\,Hz；随后四个物理步使用该周期的目标。模型选择一次 0.5\,s 踢球技能，意味着底层约运行 25 次 ONNX 推理、100 个物理步，不是 VLM 需要输出 100 次关节命令。

三个 ONNX 会话在 CPU 上运行：一个行走策略复用零速度、前后左右及转向命令，另两个是左右踢球。切换技能时保留上一时刻 14 维动作历史；不会每切换一次都假设机器人刚复位。导入时检查模型输入输出尺寸、关节顺序、HOME、动作缩放与模型哈希，防止同名模型接口不兼容。

\begin{longtable}{P{4.0cm}P{10.6cm}}
\hline
61 维底层输入 & 内容\\
\hline
\endhead
0--2、3--5 & 躯干角速度、躯干坐标中的单位重力方向，各 3 维。\\
6--19、20--33 & 14 个关节相对 HOME 的角度、14 个关节角速度。\\
34--47 & 上一次 14 维网络原始动作。\\
48--50 & 前向、横向和偏航速度命令；踢球调用使用零 twist。\\
51--54、55--60 & 头部 4 维、身体 6 维附加命令，本实验调用保持为零。\\
\hline
\end{longtable}

ONNX 输出形状为 \texttt{float32[1,14]}，按 \texttt{HOME + raw\_action} 得到关节目标，单位弧度，再交给执行器模型。当前仿真路径记录目标超出关节范围的幅度，未把所有目标自动裁剪成实机安全指令；不能把仿真通过等同于实体可直接执行。

VLM 则是低频、逐步决策：一次看图和生成完成，执行短动作并恢复，再请求下一次；没有固定 50\,Hz 的视觉语言推理频率。模型等待期间，实时物理线程继续执行 stand 平衡。离线规则采集不需等待大模型，可比墙钟实时更快；在线 VLM 的任务相对仿真时间会包括推理期间的站立时间。这是训练示范与部署之间仍要检验的时间分布差异。

本节核对：\path{runtime.py}:142--206、244--279；\path{skills.py}:25--92；本地 \path{microduck_deploy/policy.py}:163--222。

\section{问答、结束、澄清和停止分别做什么}

\begin{itemize}
\item \textbf{answer：} 当前任务是“看到了什么”，应根据当前图像回答，不先走路找更多东西。训练中的描述主要由规则球检测生成，语言覆盖集中在球颜色与方位。
\item \textbf{finish：} 模型报告任务已经结束。一次前进指令执行一个短脉冲并收到完成反馈后即可 finish；踢球则要继续处理搜索、接近等多个步骤。模型自述完成不等于独立物理验收成功。
\item \textbf{clarify：} 用户目标不明确时请求澄清，例如只说“踢球”却无法确定颜色。已经明确要求踢红球但红球暂时不在画面中，应先尝试寻找，不能把自己的搜索子任务推回给用户。
\item \textbf{stop：} 操作员说“停下”通过独立停止处理，不需要等待 VLM 生成一种 stop 标签；执行器取消当前脉冲并继续站立。
\end{itemize}

非动作响应均需 \texttt{skill=null, seconds=0}。没有独立目标验收时，程序对 finish 的结果标作 \texttt{model\_reported}，并提示尚未独立验收；问答为 \texttt{answered}，澄清为 \texttt{clarification}。带独立测试目标时，还会检查只发生一次有效踢球并通过整段轨迹验收，才标作 success。

默认单任务最多 40 次 VLM 调用，允许配置到 120 次；单次推理默认超时 45\,s，另有任务总墙钟 600\,s 上限。数据审计显示 457 个成功踢球示范的决策次数中位数为 21、95 分位约 56.8、最大 117，40 个超过默认 40 次。因此完整闭环评测宜显式用 \texttt{--vlm-max-steps 120}，并分别报告次数耗尽、推理超时和物理失败；最长示范距离 120 上限也只有 3 次余量。

\section{当前运行服务与正在训练的适配器}

现有 \texttt{serve\_vlm.sh} 启动的是本地 llama.cpp 服务，读取官方 Qwen3-VL-2B-Instruct GGUF：语言权重文件为 Q4\_K\_M，视觉投影相关文件为 Q8\_0。它通过回环地址提供兼容 chat-completions 的接口，不需要在每次任务时调用付费远程 API。

正在训练的 QLoRA 使用另一条 Hugging Face／Transformers 路径，保存的是适配器权重。训练目录出现 checkpoint 不会自动改变 llama.cpp 正在读取的 GGUF；当前启动代码没有自动加载这个 LoRA。后续需要明确完成合并／转换或新增兼容的推理服务，再把同一输入协议接入 \texttt{SkillTaskSession}，比较基础模型与适配器模型的同场景表现。

这一连接完成前，本报告可以证明数据和训练进程存在、训练损失在计算、模型接口与模拟执行链可用；不能把这些事实写成“微调模型已经在仿真中自动完成红黄蓝踢球”。更不能推出实体部署已经通过。本轮报告的运行范围仍是 HD1910 仿真。

\begin{notebox}{怎样理解当前系统名称}
第一阶段是预训练 VLM 驱动已有技能；微调的学习目标是视觉语言到技能及参数的映射，可按动作层级称为分层视觉语言动作策略。技术描述应同时写清：输入为双 RGB、原指令、必要状态与历史；输出为五字段技能 JSON；上层 LoRA 可训练，底层 ONNX 固定。这个描述比只争论“叫 VLM 还是 VLA”更能说明系统实际上做了什么。
\end{notebox}

本章核对：\path{skill_tasks.py}:17--37、185--315；\path{skill_vlm.py}:297--402；\path{setup_vlm.py}:30--33、280--308；\path{demo.py}:188--200。源码路径、样本来源及统计证据见同目录 \path{sources_local_data.json}。
